\documentclass[conference]{IEEEtran}
% Overleaf runs BibTeX from the project top level even when the selected main
% document lives in paper/.  Resolve one prefix for both a full-repository
% Overleaf project and a standalone upload of the paper/ directory.
\ifdefined\PaperRoot\else
  \IfFileExists{paper/references.bib}{%
    \newcommand{\PaperRoot}{paper}%
  }{%
    \newcommand{\PaperRoot}{.}%
  }%
\fi
% \IEEEoverridecommandlockouts is intentionally disabled: there is no funding
% footnote in the anonymous submission.
% IEEE conference template version as of 6/27/2024.

\usepackage{cite}
\usepackage{amsmath,amssymb,amsfonts}
\usepackage{array}
\usepackage{booktabs}
\usepackage{graphicx}
\usepackage{microtype}
\usepackage{tabularx}
\usepackage{textcomp}
\usepackage{url}
\usepackage{xcolor}
% Keep IEEE's typography and page geometry while avoiding oversized gaps
% around floats.
\setlength{\textfloatsep}{7pt plus 2pt minus 2pt}
\setlength{\floatsep}{6pt plus 2pt minus 2pt}
\setlength{\dbltextfloatsep}{7pt plus 2pt minus 2pt}
\setlength{\dblfloatsep}{6pt plus 2pt minus 2pt}
\def\BibTeX{{\rm B\kern-.05em{\sc i\kern-.025em b}\kern-.08em
    T\kern-.1667em\lower.7ex\hbox{E}\kern-.125emX}}

\title{Gradient Descent Optimization of Plucked-String Musical Phrases via
Cumulative Energy Loss}
\author{\IEEEauthorblockN{Anonymous Authors}
\IEEEauthorblockA{Institution withheld for double-blind review\\
contact@example.invalid}}

% LaTeX control words cannot contain underscores, so the requested
% \l_one and \l_two source names are represented by \lone and \ltwo.
\newcommand{\lone}{\ensuremath{\mathcal{L}_1}}
\newcommand{\ltwo}{\ensuremath{\mathcal{L}_2}}
\newcommand{\ddsp}{DDSP}
\newcommand{\mss}{MSS}
\newcommand{\sot}{SOT}
\newcommand{\tfw}{\ensuremath{\mathrm{TF}\mathcal{W}_2}}
\newcommand{\bicul}{\ensuremath{\mathrm{BiCu}\mathcal{L}}}
\newcommand{\redtodo}[1]{\textcolor{red}{#1}}

\begin{document}
\maketitle

\begin{abstract}
Differentiable digital signal processing enables interpretable, label-free
analysis by synthesis, yet fundamental frequency and onset remain difficult to
recover: changing either displaces coherent energy, while waveform and
pointwise spectral losses do not encode the displacement direction. External
detectors bypass this geometry but prevent the reconstruction objective from
correcting their estimates. We instead optimise a sparse sequence of plucks,
each defined by onset and fundamental frequency, through a differentiable
digital waveguide whose state and feedback delay expose onset-placement and
fractional-delay effects. Our bidirectional cumulative-energy loss compares
normalised spectral-power accumulations in both directions of time and
frequency. Across 6000 fits of one- to eight-pluck phrases, it recovers onset
and pitch more reliably than waveform, spectral, framewise-transport, and joint
time--frequency-transport objectives, and is the only tested loss with
target-directed descent at every sampled point in the single-pluck landscape.
Crossing four onset-placement and four delay-interpolation methods further
identifies Fourier placement with causal Thiran propagation as the most
consistently target-directed combination. These results support direct
audio-supervised optimisation of pitch, event timing, and delay-line length.
\end{abstract}

\begin{IEEEkeywords}
differentiable digital signal processing, digital waveguide, audio loss,
analysis by synthesis, event parameter estimation
\end{IEEEkeywords}

\section{Introduction}
Differentiable digital signal processing (\ddsp) places interpretable audio
processors in automatic-differentiation graphs for analysis by synthesis
\cite{engel2020ddsp,caspe2022ddx7}. Sound matching usually represents control
as either a global vector for a fixed or approximately linear time-invariant
processor, or dense frame-rate trajectories. Dense control suits continuously
driven instruments and evolving modulation, but is excessive when musical
changes occur at a few discrete instants.

Music Information Retrieval (MIR), by contrast, represents performances as
note events with onset, pitch, duration, and velocity
\cite{hawthorne2018onsets}. For plucked and struck instruments, we use
\begin{equation}
    \mathcal{E}=\{(t_i,f_{0,i})\}_{i=1}^{C},
\end{equation}
where pluck $i$ has onset $t_i$ and fundamental frequency $f_{0,i}$; other
controls may still vary. This gives audio-domain sound matching a note-centred
representation without a dense trajectory. Tablas de Paula et al.\ likewise
use an event-prediction encoder with a differentiable Karplus--Strong
synthesizer \cite{tablasdepaula2026karplus}, but find unreliable joint onset
and pitch gradients with standard audio losses, while synthetic parameter
supervision transfers poorly to recordings. Event parameterisation alone is
therefore insufficient.

Fundamental frequency and onset instead displace coherent energy. Pitch
perturbations alter phase throughout a note, making waveform-distance surfaces
oscillatory; pointwise spectral distances compare only corresponding bins and
become weak or misleading when candidate and target energy do not overlap
\cite{turian2020loss,hayes2023sinusoidal}. Many DDSP systems bypass this
geometry with external $f_0$ estimates or synthetic parameter labels
\cite{engel2020inverse}, respectively preventing audio-objective correction or
risking a simulation-to-recording gap.

Optimal transport addresses displacement directly: spectral optimal transport
(\sot) moves energy along frequency within each frame \cite{torres2024sot},
whereas joint time--frequency optimal transport (\tfw) also moves it across
time \cite{fabiani2024jtfot}, subject to the metric and relative axis scale. We
instead introduce the bidirectional cumulative-energy loss (\bicul), comparing
normalised spectral-power accumulations in both directions of both axes without
detectors or target-event coordinates.

We evaluate direct event optimisation through a differentiable digital
waveguide: excitation time sets onset, propagation delay sets $f_0$, and
recursive state couples plucks. This compact model yields a rich, decaying
string signal \cite{valimaki1996plucked,depaula2026waveguides} while exposing
the renderer Jacobian: sub-sample onset placement and fractional-delay
interpolation can give similar audio but different derivatives
\cite{tablasdepaula2026karplus}. We contribute the event formulation, \bicul's
gradient and phrase-recovery evaluation, and a controlled study of how onset
placement, delay interpolation, and recursive state affect those gradients.

\section{Cumulative energy losses}
\label{sec:losses}
Let $x$ be the target, $\theta$ the candidate-event coordinates, and
$\hat x=R_\rho(\theta)$ the output of renderer realisation $\rho$. Although the
loss is synthesizer-independent, its parameter gradient contains the renderer
Jacobian:
\begin{equation}
\nabla_\theta\mathcal L(x,R_\rho(\theta))=
J_{R_\rho}(\theta)^\mathsf{T}\nabla_{\hat x}\mathcal L(x,\hat x).
\label{eq:pullback}
\end{equation}
We therefore compare losses under one fixed renderer, then fix \bicul{} to
compare renderer realisations.

\subsection{Reference objectives}
\lone{} is mean absolute sample error and \ltwo{} is mean squared sample
error. Following Schw\"ar and M\"uller \cite{schwar2023mss}, \mss{} sums
squared \ltwo{} distances between $\log(1+|\operatorname{STFT}|)$ magnitudes.
Periodic prime-length flat-top windows
$\{17,31,67,127,257,509\}$ with hops $\{8,15,33,63,128,254\}$ suppress
spectral-leakage sidelobes without aligning DFT bins across resolutions.

Torres et al.'s \sot{} \cite{torres2024sot} combines frame-normalised
transport on FFT-512/hop-64 flat-top STFT power, with their unbalanced quantile
cutoff, and $0.05$ times Hann-windowed linear-magnitude MSS. Fabiani et al.'s
\tfw{} \cite{fabiani2024jtfot} globally normalises a periodic-Hann
FFT-256/hop-128 magnitude STFT and projects it along ten fixed angles;
frequency is in Hz and $\widetilde t=1000t$ in Hz, equating 1~s with 1000~Hz.

\subsection{Bidirectional cumulative energy}
Let $q=(d_t,d_f)$ select time and frequency accumulation directions, and let
$P,\hat P$ be periodic-Hann FFT-256/hop-64 STFT power:
\begin{equation}
\begin{aligned}
S_q&=C^f_{d_f}(C^t_{d_t}(P)),\\
\hat S_q&=C^f_{d_f}(C^t_{d_t}(\hat P)),\\
m_q&=\max S_q,\qquad R_q=\sqrt{\max(S_q/m_q,10^{-12})},\\
\hat R_q&=\sqrt{\max(\hat S_q/m_q,10^{-12})},\\
\ell_q&=\operatorname{RMS}(R_q-\hat R_q).
\end{aligned}
\end{equation}
\bicul{} averages $\ell_q$ over
$d_t\in\{\rightarrow,\leftarrow\}$ and
$d_f\in\{\uparrow,\downarrow\}$. Displacement changes cumulative values from
its location to a boundary, exposing whether energy is early/late or low/high;
both orientations avoid favouring one boundary. The construction uses no event
coordinates, detectors, or salience weights.

\subsection{Conditional loss geometry}
Figure~\ref{fig:landscapes} fixes the reference renderer from
Sec.~\ref{sec:renderer} and evaluates a 160-Hz, 1.0-s event on an exact
25-by-25 lattice spanning $\pm0.8$~s linearly and $\pm1$ octave
logarithmically, including the centred target. Surfaces are independently
min--max normalised; white arrows are autograd unit-descent directions on an
11-by-11 sublattice, excluding the target. A direction is target-directed when
its angle to the target displacement is below $90^\circ$. Of 120 non-target
directions, \lone{}, \ltwo{}, \mss{}, \sot{}, \tfw{}, and \bicul{} achieve
60.0\%, 49.2\%, 42.5\%, 51.7\%, 75.8\%, and 100.0\%, respectively.

\begin{figure}[t]
\centering
\setlength{\abovecaptionskip}{0pt}
\IfFileExists{\PaperRoot/figures/loss_landscapes.pdf}{
  \includegraphics[width=\columnwidth]{\PaperRoot/figures/loss_landscapes.pdf}
}{\fbox{\parbox[c][58mm][c]{0.94\columnwidth}{\centering
  Loss-direction figure pending checksum-complete generation.}}}
\vspace{-5.8mm}
\caption{Normalised single-pluck loss landscapes; the cross marks the target
and arrows show unit descent.}
\label{fig:landscapes}
\end{figure}

\section{Differentiable phrase renderer}
\label{sec:renderer}
Figure~\ref{fig:waveguide} fixes the physical topology of $R_\rho$: two
oppositely travelling rails, each divided by the pluck into nut and bridge
sections. Their summed length determines $f_0$ and their ratio the pluck
position \cite{depaula2026waveguides}, fixed here to 0.23 from the bridge. The
nut gain is $-1$ and the bridge reflection is $-H_L(z)$, where
\begin{equation}
H_L(z)=g\frac{1-a_1}{1-a_1z^{-1}},\qquad g=0.99,\quad a_1=0.2.
\end{equation}
$g$ sets low-frequency round-trip gain and $a_1$ the increasing high-frequency
attenuation, producing a complex decay whose RT60 varies with $f_0$ as in
Fig.~3 of \cite{tablasdepaula2026karplus}. Output is bridge-incident travelling
velocity.

\begin{figure}[t]
\centering
\includegraphics[width=\columnwidth]{\PaperRoot/figures/Waveguide.pdf}
\caption{Two-rail digital waveguide renderer.}
\label{fig:waveguide}
\end{figure}

\subsection{Differentiable onset placement and interpolation}
\begin{figure*}[t!]
\centering
\setlength{\abovecaptionskip}{0pt}
\IfFileExists{\PaperRoot/figures/state_realization.pdf}{
  \includegraphics[width=\textwidth]
  {\PaperRoot/figures/state_realization.pdf}
}{\fbox{\parbox[c][95mm][c]{0.96\textwidth}{\centering
  State-realisation comparison pending checksum-complete generation.}}}
\vspace{-7.0mm}
\caption{Reset-DF2 target spectrogram and middle-event \bicul{} sweeps for
collapsed DF2 and Literal waveguides.}
\label{fig:state-realization}
\end{figure*}
Audio is float64, 2~s long, and sampled at 4~kHz. Each event excites the
zero-onset half-raised-cosine (HRC)
\begin{equation}
e_0(t)=\begin{cases}
\dfrac{A_e}{2}[1-\cos(\pi t/T_e)],&0\leq t<T_e,\\
0,&\text{otherwise,}
\end{cases}
\end{equation}
with $A_e=0.8$ and $T_e=10$~ms. Four methods place it continuously: Naive
samples the analytic HRC at the requested onset; Lagrange applies a six-tap,
fifth-order maximally-flat FIR shift; Fourier applies $e^{-j\omega D}$ using
FLAMO \cite{dalsanto2025flamo}; and Thiran applies a causal first-order
all-pass shift.

For propagation $\mathcal D_p(D;z)\approx z^{-D}$, Linear uses two-tap order-1
Lagrange sections; Lagrange uses order 5 on the long nut section and order 1 on
the short bridge section; Fourier samples the delay response; and Thiran uses
all-pass orders 3 and 1 on those respective sections. Lagrange and Thiran are
maximally-flat FIR and all-pass designs \cite{laakso1996splitting}; both use
order 1 on the bridge side because its delay can be too short for the
higher-order form. Rails are tuned with the realised fractional-delay and loop-
filter phase.

Table~\ref{tab:interpolation-onset} is the complementary slice to
Fig.~\ref{fig:landscapes}: it fixes \bicul{} while target and candidate share
each onset/propagation method on the same lattice, isolating derivative geometry
from renderer mismatch. The Fourier waveguide is a FLAMO module whose
65,536-point transform is 8.192 times the 8000-sample horizon to suppress time
aliasing; real DC/Nyquist bins and conjugate negative frequencies enforce
Hermitian symmetry.

\begin{table}[h]
\centering
\caption{\normalfont Target-directed \bicul{} descent directions (\%) by onset
placement (rows) and propagation interpolation (columns); higher is better.}
\label{tab:interpolation-onset}
\resizebox{\columnwidth}{!}{%
\begin{tabular}{lrrrr}
\toprule
$\downarrow\,t\;\backslash\;f_0\,\rightarrow$
  & Linear & Lagrange & Fourier & Thiran \\
\midrule
Naive    & 52.5\% & 62.5\% & 68.3\% & 68.3\% \\
Lagrange & 49.2\% & 62.5\% & 70.0\% & 63.3\% \\
Fourier  & 66.7\% & 78.3\% & 88.3\% & 100.0\% \\
Thiran   & 65.8\% & 79.2\% & 99.2\% & 99.2\% \\
\bottomrule
\end{tabular}
}
\end{table}

Table~\ref{tab:interpolation-onset} identifies Fourier onset placement with
Thiran propagation as the only 120/120 result. We therefore adopt this
Fourier-placement/causal-Thiran pairing for Fig.~\ref{fig:landscapes} and all subsequent
experiments.

\subsection{Phase interference and string reset}
Each constant event regime collapses to an exact rational transfer, evaluated
through PhilTorch's direct-form-II (DF2) interface and TorchLPC backend
\cite{yu2025philtorch}. Equivalent fixed-coefficient filter rearrangements can
differ after a coefficient change \cite{yu2024torchlpc}: collapsed DF2
compresses rails, reflections, and bridge termination into one history that new
coefficients reinterpret. Our Literal control instead retains separate physical
rail and bridge-filter histories. They are equivalent within a constant regime
but need not agree across coefficient changes.

Figure~\ref{fig:state-realization} sweeps the middle event in a three-pluck
target $(110,0.5),(160,1.0),(280,1.5)$ in Hz and seconds. Persistent DF2 and
Literal slopes differ, but both are rugged and repeatedly misdirected. Clearing
all filter and delay histories at each event boundary makes them agree and
restores smooth, target-directed slopes, so subsequent experiments use reset.
We hypothesise that phase interference between carried oscillation and the new
pluck corrugates persistent loss; future work should isolate this mechanism.

\section{Experiments}
\subsection{Targets}
Seed 2027 generates 150 scrambled-Latin-hypercube targets at each of one, two,
four, six, and eight events. Pitches are log-mapped to 80--320~Hz; onsets span
0.2--1.8~s with at least 50~ms separation.

\subsection{Optimisation}
All candidate pitches start at 160~Hz and onsets at equal-cell centres over
0.2--1.8~s; independent logits bound coordinates to those same onset and pitch
ranges. Adam uses learning rate 0.05, default moments and epsilon, and no weight
decay. Dividing each loss by its initial value equalises numerical scale without
changing reported values. A relative raw-loss decrease of $10^{-4}$ resets
patience. At patience 100 and 200, we restore the strict best, multiply the
learning rate by 0.3 (floor $10^{-5}$), and clear optimiser state. We stop when
patience reaches 250 or after 3000 updates and retain the best iterate.

\subsection{Loss variants}
\label{sec:loss-variants}
The $1\,\mathrm{s}=1\,\mathrm{oct}$ variant changes only \tfw{} coordinates to
$\log_2[\max(f,20)/20]$ octaves and seconds, equating one second and one octave.
Log-Q changes only \bicul's final RMS, replacing uniform STFT-bin weights with
quadrature in seconds and log-frequency octaves. For
$E_{d_t,d_f}=\hat R_q-R_q$,
\begin{equation}
\ell^{\rm quad}_{d_t,d_f}=\left[\frac{\sum_{k,n}
w^f_{d_f,k}w^t_{d_t,n}E_{d_t,d_f}(k,n)^2}
{\sum_{k,n}w^f_{d_f,k}w^t_{d_t,n}}\right]^{1/2},
\end{equation}
where each $w$ is its direction-aligned adjacent-coordinate interval; all other
operations are unchanged. With Sec.~\ref{sec:losses}'s six objectives, these
variants yield 6000 fits: 150 targets at five event counts across 40 loss--count
conditions.

\subsection{Evaluation}
Because predictions are unordered, Hungarian assignment pairs them one-to-one
with targets by minimising
\begin{equation}
d_{ij}=\left[\frac{\log_2(\hat f_i/f_j)}{2}\right]^2+
\left[\frac{\hat t_i-t_j}{1.6}\right]^2.
\end{equation}

\begin{table*}[!t]
\centering
\caption{\normalfont Median Hungarian-matched pitch and onset errors over 150
phrases per event count; lower is better.}
\label{tab:pilot-compact}
\begingroup
\footnotesize
\setlength{\tabcolsep}{2.5pt}
\renewcommand{\arraystretch}{1.02}
\begin{tabularx}{\columnwidth}{@{}l*{5}{>{\centering\arraybackslash}X}@{}}
\toprule
Loss & 1 Event & 2 Events & 4 Events & 6 Events & 8 Events \\
\midrule
& \multicolumn{5}{c}{$\Delta f_0$ (Cents)} \\
\lossLone{} & 948 & 855 & 732 & 624 & 552 \\
\lossLtwo{} & 921 & 894 & 800 & 687 & 600 \\
\lossMSS{} & \textit{96.8} & \textit{262} & \textit{316} & \textit{278} & \textit{295} \\
\lossSOT{} & 840 & 837 & 877 & 800 & 757 \\
\lossTFW{} & \underline{0.173} & \underline{4.27} & \underline{98.4} & \underline{187} & \underline{180} \\
\lossBiCum{} & {\bfseries\boldmath $0.000802$} & {\bfseries\boldmath $0.00139$} & {\bfseries\boldmath $38.4$} & {\bfseries\boldmath $145$} & {\bfseries\boldmath $152$} \\
\addlinespace[2pt]
\midrule
\addlinespace[1pt]
& \multicolumn{5}{c}{$\Delta t$ (ms)} \\
\lossLone{} & \textit{65.8} & \textit{190} & 235 & 186 & 185 \\
\lossLtwo{} & 417 & 378 & 242 & 214 & 191 \\
\lossMSS{} & 210 & 219 & \textit{201} & \textit{155} & \textit{154} \\
\lossSOT{} & 97.1 & 228 & 206 & 192 & 167 \\
\lossTFW{} & \underline{0.257} & \underline{6.46} & \underline{85.7} & \underline{106} & \underline{133} \\
\lossBiCum{} & {\bfseries\boldmath $0.000321$} & {\bfseries\boldmath $0.000345$} & {\bfseries\boldmath $31.4$} & {\bfseries\boldmath $93.2$} & {\bfseries\boldmath $100$} \\
\bottomrule
\end{tabularx}
\endgroup

\end{table*}

We report the across-phrase median of within-phrase mean absolute pitch cents
and onset milliseconds. Periodic-Hann FFT-256/hop-64 log-spectral distance
(LSD) gives an objective-independent audio metric for every phrase.

\section{Results}
Table~\ref{tab:pilot-compact} reports all 40 conditions. Bold, underlined, and
italic entries are the three lowest medians per metric and event count; hooked
rows denote Sec.~\ref{sec:loss-variants}'s variants. Figure~\ref{fig:lsd} shows
all LSD values with white interquartile bars and outlined medians.

\begin{figure}[!t]
\centering
\setlength{\abovecaptionskip}{0pt}
\IfFileExists{\PaperRoot/figures/lsd_violin.pdf}{
  \includegraphics[width=\columnwidth]{\PaperRoot/figures/lsd_violin.pdf}
}{\fbox{\parbox[c][45mm][c]{0.96\textwidth}{\centering
  Figure 4 LSD distributions pending checksum-complete regeneration.}}}
\vspace{-8.0mm}
\caption{Per-phrase LSD distributions over the same targets; lower is better.}
\label{fig:lsd}
\end{figure}

\newpage
\section{Conclusion}
We introduced direct audio recovery of sparse pluck onsets and pitches and
\bicul{}, a bidirectional cumulative-STFT-power loss. Across 6000 fits spanning
one to eight events, \bicul{} was fully target-directed on the sampled
single-pluck landscape and recovered multi-pluck timing and pitch more reliably
than classical, spectral, and optimal-transport objectives; Log-Q further
improved descriptive multi-event results. Our renderer study identified Fourier
onset placement with Thiran propagation as the strongest tested pair and found
that event-wise reset corrected poor gradients in both collapsed DF2 and
Literal persistent waveguides.

Future work should address expressive techniques including vibrato, higher
sample rates, state treatment without event-wise reset, neural-network
training, and real-world recordings.

\clearpage
\bibliographystyle{IEEEtran}
\typeout{ICASSP bibliography database: \PaperRoot/references.bib}
\bibliography{\PaperRoot/references}
\end{document}
